\section{Optimisation et Levée des Verrous sur AMD}

\begin{frame}{Stratégie d'optimisation : Briser l'aliasing}
    \begin{itemize}
        \fontsize{9}{10}\selectfont
        \item L'analyse sur architecture AMD a démontré que l'ambiguïté mémoire bloquait la vectorisation automatique.
        \item \textbf{Objectif} : Fournir des indices sémantiques au compilateur pour forcer l'usage du SIMD.
    \end{itemize}
    \begin{block}{Leviers techniques implémentés}
        \fontsize{8}{9}\selectfont
        \begin{itemize}
            \item \textbf{\_\_restrict} : Garantit l'absence de chevauchement entre les pointeurs de tableaux.
            \item \textbf{\#pragma omp simd} : Force explicitement la vectorisation de la boucle interne.
            \item \textbf{Clause collapse(2)} : Fusionne les boucles imbriquées pour augmenter l'efficacité.
        \end{itemize}
    \end{block}
\end{frame}

\begin{frame}[fragile]{Transformation de la routine selon CQA}
    \fontsize{6}{7}\selectfont
    \begin{columns}[T]
        \begin{column}{0.48\textwidth}
            \begin{alertblock}{Version Originale}
\begin{verbatim}
void resid(double* u, double* v, 
           double* r, ...) {
  for (int i = 0; i < n; i++) {
    for (int j = 0; j < m; j++) {
      r[i][j] = v[i][j] - u[i][j]...;
    }
  }
}
\end{verbatim}
            \end{alertblock}
            \centering \textit{Le compilateur suspecte un risque d'aliasing et refuse le SIMD.}
        \end{column}
        
        \begin{column}{0.48\textwidth}
            \begin{exampleblock}{Version Optimisée}
\begin{verbatim}
// Code vectoriel (AVX/SIMD)
void resid(double* __restrict u, 
           double* __restrict v, ...) {
  #pragma omp simd collapse(2)
  for (int i = 0; i < n; i++) {
    for (int j = 0; j < m; j++) {
      r[i][j] = v[i][j] - u[i][j]...;
    }
  }
}
\end{verbatim}
            \end{exampleblock}
            \centering \textit{La sémantique \_\_restrict lève le verrou et active la vectorisation.}
        \end{column}
    \end{columns}
\end{frame}

\begin{frame}{Impact Global sur AMD : Analyse des routines}
    \inputImage{figures/top_functions_amd_r1_opt.png}{Routines AMD optimisées (Analyse -R1)}{0.95\textwidth}{!}{fig:global_opt_amd}
    
    \begin{table}
        \fontsize{6}{7}\selectfont
        \centering
        \begin{tabular}{|l|c|c|}
            \hline
            \textbf{Fonction} & \textbf{\% Temps CPU (Initial)} & \textbf{\% Temps CPU (Optimisé)} \\ \hline
            \texttt{resid}    & 47,85 \%            & \textbf{29,57 \%}    \\ \hline
            \texttt{zran3}    & -                               & \textbf{25,95 \%}    \\ \hline
            \texttt{psinv}    & 19,36 \%            & \textbf{11,68 \%}    \\ \hline
            \texttt{interp}   & 9,50 \%             & \textbf{5,42 \%}     \\ \hline
        \end{tabular}
    \end{table}
\end{frame}

\begin{frame}{Analyse de l'anomalie de détection sur MAQAO}
    \begin{block}{Observation du paradoxe}
        \fontsize{8}{9}\selectfont
        \begin{itemize}
            \item \textbf{Performance réelle} : On observe une diminution nette du temps d'exécution relatif (ex: \texttt{resid} passe de 47,85 \% à 29,57 \%), validant l'efficacité des optimisations.
            \item \textbf{Score réelle} : Le score de compilation reste figé à \textbf{16,7 \%} malgré l'injection de sémantique vectorielle.
        \end{itemize}
    \end{block}

    \begin{block}{Hypothèses sur l'origine du blocage}
        \fontsize{8}{9}\selectfont
        \begin{itemize}
            \item \textbf{Visibilité des flags} : Difficulté de MAQAO à extraire les drapeaux d'optimisation des métadonnées du binaire C++.
            \item \textbf{Liaison Makefile} : Un potentiel défaut de "binding" des flags d'optimisation lors de la phase finale d'édition de liens.
            \item \textbf{Metadata Compilateur} : Le compilateur exécute les optimisations mais ne les consigne pas de manière lisible pour l'outil d'analyse.
        \end{itemize}
    \end{block}
\end{frame}